% ch13 译文：原文附录 A §A.a（main.tex 第 4656--5103 行）
\appendix

\section{\v{C}ech 复形的 Eilenberg--Zilber 收缩}\label{sec:Eilenberg-Zilber-constraction}

\subsection{链复形的 Eilenberg--Zilber 收缩}

本小节沿用 \cite[\S 3]{Franz} 中对 Eilenberg--Mac Lane 古典构造
\cite{EilenbergMacLane1953,EilenbergMacLane1954} 的表述，为与有序集相伴的链复形
构造 \emph{Eilenberg--Zilber 映射}，并建立这些映射之间的关系。

设 $I$ 为有序集。与 $I$ 相伴的链复形
$
C_{\bullet}(I)
:=
\left(\bigoplus_{p\in\mathbb{Z}}C_p(I),\partial\right)
$
的各项为
$$
C_p(I)
:=
\mathbb{Z}\left<\left\{\left.(i_0,\dots,i_p)\in I^{p+1}\right|i_0<\cdots<i_p\right\}\right>
\quad\text{对 }p\ge0,
$$
当 $p<0$ 时 $C_p(I):=0$；边界算子为
\begin{align*}
\partial:C_p(I)
&\rightarrow
C_{p-1}(I)
\qquad
(p\ge 1),
\\
(i_0,\dots,i_p)
&\mapsto
\sum_{\nu=0}^p
(-1)^{\nu}
(i_0,\dots,\hat{i_{\nu}},\dots,i_p).
\end{align*}
对 $I$ 中满足 $i_0<\cdots<i_p$ 的指标以及 $\sigma\in S_{p+1}$，记
$
(i_{\sigma(0)},\dots,i_{\sigma(p)})
:=
\operatorname{sign}(\sigma)
(i_0,\dots,i_p)
$；
若 $i_0,\dots,i_p$ 中有两个相同，则记 $(i_0,\dots,i_p):=0$。

给定两个有序集 $I$ 与 $J$，考察两个自然的链复形
$$
C_{\bullet}(I\times J)
\quad\text{与}\quad
C_{\bullet}(I)\otimes C_{\bullet}(J).
$$
前者是与带字典序的 $I\times J$ 相伴的链复形；后者是链复形 $C_{\bullet}(I)$ 与
$C_{\bullet}(J)$ 的通常张量积，其各次分量为
$$
(C_{\bullet}(I)\otimes C_{\bullet}(J))_t
:=
\bigoplus_{p+r=t}
\left(C_p(I)\otimes C_r(J)\right)
$$
边界算子为
\begin{align*}
\partial:
(C_{\bullet}(I)\otimes \partial_{\bullet}(J))_t
&\rightarrow
(C_{\bullet}(I)\otimes \partial_{\bullet}(J))_{t-1}
\\[2mm]
(i_0,\dots,i_p)\otimes(j_0,\dots,j_r)
&\mapsto
\partial((i_0,\dots,i_p))\otimes(j_0,\dots,j_r)
+
(-1)^p
(i_0,\dots,i_p)\otimes\partial((j_0,\dots,j_r)).
\end{align*}

\begin{yizhu}[分次符号与一处下标笔误]
式子 $C_{\bullet}(I)\otimes\partial_{\bullet}(J)$ 里的 $\partial_{\bullet}(J)$ 是原文笔误，
应为 $C_{\bullet}(J)$：这个复形的第 $t$ 次分量是 $\bigoplus_{p+r=t}C_p(I)\otimes C_r(J)$，
并非把 $C_{\bullet}(I)$ 与 $\partial_{\bullet}(J)$ 张量起来。译文照原文排印，读者按
$C_{\bullet}(J)$ 理解即可。另一方面，$(-1)^p$ 不是笔误：把 $\partial$ 推过第一个因子是分次
交换，按 Koszul 法则要付出 $\deg(i_0,\dots,i_p)=p$ 个负号；后面每个公式的符号都从这条规则来，
符号错一位结果就全错。
\end{yizhu}

% [原文 p.48]
我们断言这两个链复形彼此同伦等价，办法是显式构造链映射与同伦，它们合称
\emph{Eilenberg--Zilber 映射}：
\begin{equation}\label{eqn:EZMaps}
\begin{tikzcd}[arrow style=tikz,>=stealth,row sep=2.5em,column sep=3.5em]
C_{\bullet}(I)\otimes C_{\bullet}(J)
  \arrow[r, shift left, "S"]
&
C_{\bullet}(I\times J)
  \arrow[l, shift left, "T"]
  \arrow[loop right, <-, "\mathbb{H}"]
\end{tikzcd}
\end{equation}

\begin{defn}\label{defn:EZMaps}
(1) \emph{洗牌映射}（shuffle map）由下式定义：
\begin{center}
\resizebox{\textwidth}{!}{$\begin{aligned}
S:
C_{\bullet}(I)\otimes C_{\bullet}(J)
&\rightarrow
C_{\bullet}(I\times J)
\\[2mm]
(i_0,\dots,i_p)\otimes(j_0,\dots,j_r)
&\mapsto
\text{
\scalebox{0.8}{
$\sum_{
0=l_0\le l_1\le\cdots\le l_{p+1}=r
}
(-1)^{l_1+\cdots+l_p}
\big(
(i_0,j_0),\dots,(i_0,j_{l_0}),
(i_1,j_{l_1}),\dots,(i_1,j_{l_2}),
\dots,
(i_p,j_{l_p}),\dots,(i_p,j_r)
\big),
$}}
\end{aligned}$}
\end{center}
其中每一步恰有一个 $i$ 或 $j$ 的下标增加 $+1$。

\noindent
(2) \emph{Alexander--Whitney 映射}由下式定义：
\begin{align*}
\begin{split}
T:
C_{\bullet}(I\times J)
&\rightarrow
C_{\bullet}(I)\otimes C_{\bullet}(J)
\\[0mm]
\big(\alpha_0=(i_0,j_0),\dots,\alpha_t=(i_t,j_t)\big)
&\mapsto
\sum_{p=0}^t
(i_0,\dots,i_p)\otimes(j_p,\dots,j_t).
\end{split}
\end{align*}
\end{defn}

直接验证可知 $S$ 与 $T$ 都是链映射，并且 $S\comp T$ 是恒等映射：
$$
S\comp \partial=\partial \comp S,
\qquad
T\comp \partial=\partial \comp T,
\qquad
S\comp T=\id_{C_{\bullet}(I)\otimes C_{\bullet}(J)}.
$$

% [原文 p.49]
\begin{prop}\label{prp:chain-homotopy-TS}
链映射 $T\bullet S$ 与 $\id_{C_{\bullet}(I\times J)}$ 彼此同伦等价。也就是说，存在链同伦
$$
\mathbb{H}:
C_{\bullet}(I\times J)
\rightarrow
C_{\bullet+1}(I\times J)
$$
连接 $T\bullet S$ 与 $\id_{C_{\bullet}(I\times J)}$，并满足
$$
\mathbb{H}\comp \partial+\partial\comp \mathbb{H}
=
T\comp S-
\id_{C_{\bullet}(I\times J)}.
$$
此外它还具有性质 $S\comp \mathbb{H}=0$。
\end{prop}

\begin{yizhu}[为什么只是同伦等价，不是同构]
两个复形的次数不一样：$C_{\bullet}(I)\otimes C_{\bullet}(J)$ 的一次项是若干 $p+r=t$ 的
张量积之和，$C_{\bullet}(I\times J)$ 的一次项是乘积有序集上长度为 $t$ 的链，二者维数通常不等，
所以不可能同构。$S\comp T=\id$ 说明 $T$ 是单射、$S$ 是满射；$T\comp S$ 与恒等只差一个同伦
$\mathbb{H}$，而 $\mathbb{H}$ 把次数抬高 $1$，于是二者在同调层面给出同一复形。条件
$S\comp\mathbb{H}=0$ 是「收缩」的附加要求（形变收缩的侧条件），后文把
$(S,T,\mathbb{H})$ 直接当成一套收缩来用。原文此处写作 $T\bullet S$，与前后文的 $T\comp S$
同义，$\bullet$ 只是 $\comp$ 的另一种排版。
\end{yizhu}

\begin{conceptbox}{Eilenberg–Zilber 收缩与 Čech 复形}
对有序集 $I$，复形 $C_{\bullet}(I)$ 以递增指标组 $(i_0<\cdots<i_p)$ 为基，这正是
Čech 复形在开覆盖的指标集上的模型；两个方向上的数据合起来就是双复形。洗牌映射
$S\colon C_{\bullet}(I)\otimes C_{\bullet}(J)\to C_{\bullet}(I\times J)$ 把一对递增指标组
「归并」成一个递增指标组，Alexander--Whitney 映射 $T$ 则按分界点把它劈开：
$T$ 是 $S$ 的同伦逆，二者相差的同伦正是 $\mathbb{H}$。三元组 $(S,T,\mathbb{H})$ 满足
$S\comp T=\id$、$T\comp S-\id=\mathbb{H}\comp\partial+\partial\comp\mathbb{H}$ 与
$S\comp\mathbb{H}=0$，称为 Eilenberg--Zilber 收缩。附录 A 用它把大复形（张量积侧或乘积侧）
上的计算归约到另一侧，从而给出 Čech 复形之间所需的同伦共轭。
\par\smallskip\noindent{\footnotesize 参考：Eilenberg--Zilber, \emph{On products of complexes},
Amer. J. Math. 75 (1953)；本文附录 A。}
\end{conceptbox}

\subsection{\v{C}ech 复形的 Eilenberg--Zilber 收缩}\label{subsec:Eilenberg-Zilber-contraction-Cech}

固定对象
$
(\mathcal{F}^{\bullet},\theta^{\bullet}),
$
$
(\mathcal{F}'^{\bullet},\theta'^{\bullet})
\in\Perf_G(X;\mathfrak{U})_{\dg}
$
与
$
(\mathcal{G}^{\bullet},\eta^{\bullet}),
$
$
(\mathcal{G}'^{\bullet},\eta'^{\bullet})
\in\Perf_H(Y;\mathfrak{V})_{\dg}.
$
考察两个上链复形
\begin{align*}
C_{\otimes}^{\bullet}
&:=
\hom_{
\Perf_G(X;\mathfrak{U})_{\dg}
\otimes
\Perf_H(Y;\mathfrak{V})_{\dg}
}^{\bullet}
\Big(
\big(
(\mathcal{F}^{\bullet},\theta^{\bullet}),
(\mathcal{G}^{\bullet},\eta^{\bullet})
\big),
\big(
(\mathcal{F}'^{\bullet},\theta'^{\bullet}),
(\mathcal{G}'^{\bullet},\eta'^{\bullet})
\big)
\Big),
\\[2mm]
C_{\times}^{\bullet}
&:=
\hom_{\Perf_{G\times H}(X\times Y;\mathfrak{U}\times\mathfrak{V})_{\dg}}^{\bullet}
\Big(
(\mathcal{F}^{\bullet},\theta^{\bullet})
\boxtimes
(\mathcal{G}^{\bullet},\eta^{\bullet}),
(\mathcal{F}'^{\bullet},\theta'^{\bullet})
\boxtimes
(\mathcal{G}'^{\bullet},\eta'^{\bullet})
\Big).
\end{align*}
我们把
\begin{align*}
\begin{split}
&
\hom_{\Perf_{G}(X;\mathfrak{U})_{\dg}}^l
\big(
(\mathcal{F}^{\bullet},\theta^{\bullet}),
(\mathcal{F}'^{\bullet},\theta'^{\bullet})
\big)
=
\bigoplus_{p+q=l}
\check{C}_{G}^p
\Big(
\mathfrak{U},
\mathcal{H}om_{\mathcal{O}_X}^q
\big(
(\mathcal{F}^{\bullet},\theta^{\bullet}),
(\mathcal{F}'^{\bullet},\theta'^{\bullet})
\big)
\Big),
\\[2mm]
&
\hom_{\Perf_{H}(Y;\mathfrak{V})_{\dg}}^m
\big(
(\mathcal{G}^{\bullet},\eta^{\bullet}),
(\mathcal{G}'^{\bullet},\eta'^{\bullet})
\big)
=
\bigoplus_{r+s=m}
\check{C}_{H}^r
\Big(
\mathfrak{V},
\mathcal{H}om_{\mathcal{O}_Y}^s
\big(
(\mathcal{G}^{\bullet},\eta^{\bullet}),
(\mathcal{G}'^{\bullet},\eta'^{\bullet})
\big)
\Big),
\\[2mm]
&
C_{\times}^n
=
\bigoplus_{t+u=n}
\check{C}_{G\times H}^t
\Big(
\mathfrak{U}\times\mathfrak{V},
\mathcal{H}om_{\mathcal{O}_{X\times Y}}^u
\big(
(\mathcal{F}^{\bullet},\theta^{\bullet})
\boxtimes
(\mathcal{G}^{\bullet},\eta^{\bullet}),
(\mathcal{F}'^{\bullet},\theta'^{\bullet})
\boxtimes
(\mathcal{G}'^{\bullet},\eta'^{\bullet})
\big)
\Big)
\end{split}
\end{align*}
中固定双次 $(p,q)$、$(r,s)$、$(t,u)$ 的齐次元素分别记为
$$
f:=
\Big(f_{i_0\cdots i_p}^i\Big)
_{\substack{
i\in\mathbb{Z}
\\
i_0<\cdots<i_p
}},
\quad
g:=
\Big(g_{j_0\cdots j_r}^j\Big)
_{\substack{
j\in\mathbb{Z}
\\
j_0<\cdots<j_r
}},
\quad
h:=
\Big(h_{\alpha_0\cdots \alpha_t}^k\Big)
_{\substack{
k\in\mathbb{Z}
\\
\alpha_0<\cdots<\alpha_t
}},
$$
这些符号按下面的等同来理解：
